\chapter{Rigorous Renormalization Group}
\label{ch:constructive-rigorous-renormalization-group}

\ConventionsEuclidean

This chapter formulates renormalization group transformations as maps on
normed spaces of interactions or probability measures.  The aim is to replace
informal terminology about flows by explicit maps, estimates, fixed points,
and coordinate decompositions.

\section{Gaussian Block Integration}

Let \(C\) be a positive covariance on a finite regulator space of fields and
suppose it decomposes as
\[
  C=C_<+\Gamma .
\]
Let \(d\mu_C\), \(d\mu_{C_<}\), and \(d\mu_\Gamma\) be the corresponding
Gaussian measures.  For an interaction \(V\), define the single-step
effective interaction
\[
  e^{-V_<(\phi)}
  =
  \int d\mu_\Gamma(\zeta)\,e^{-V(\phi+\zeta)} .
\]
This identity is exact at finite regulator.  A scale transformation \(S_L\)
then gives an RG map
\[
  \mathcal R_L(V)=S_L V_< .
\]
The definition contains three inputs: covariance decomposition, integration
domain, and rescaling convention.

\section{Banach Space Of Interactions}

Let \(\mathcal B\) be a Banach space of local functionals with norm
\(\|\cdot\|_{\mathcal B}\).  A typical constructive choice separates relevant
local coordinates from a polymer remainder:
\[
  V(\phi)=\sum_i g_i\mathcal O_i(\phi)+K(\phi),
\]
where \(g_i\in\R\) or \(\C\) and \(K\) is measured by a norm that includes
large-field weights and spatial decay.  The RG map is controlled on an open
set \(U\subset\mathcal B\) if
\[
  \mathcal R_L:U\longrightarrow\mathcal B
\]
is differentiable and its derivative satisfies explicit bounds in the chosen
norm.

\section{Linearization And Coordinates}

Let \(V_\ast\) be a fixed point:
\[
  \mathcal R_L(V_\ast)=V_\ast .
\]
The linearized map is
\[
  D\mathcal R_L|_{V_\ast}:\mathcal B\to\mathcal B .
\]
An eigenvector \(e_i\) with eigenvalue \(L^{y_i}\) is relevant, marginal, or
irrelevant according as \(y_i>0\), \(y_i=0\), or \(y_i<0\).  These words refer
to the declared linearized map and norm.  A coordinate \(g_i\) is a scaling
coordinate only after nonlinear terms have been put into a normal form or
controlled by a stable-manifold argument.

\begin{theorem}[Local stable manifold for an RG map]
\label{thm:rg-local-stable-manifold}
Let \(\mathcal R\) be a \(C^1\) map on a Banach space \(\mathcal B\) with
fixed point \(V_\ast\).  Suppose
\[
  \mathcal B=E_u\oplus E_s
\]
where \(E_u\) is finite dimensional, \(E_s\) is closed, and the derivative
\(D\mathcal R|_{V_\ast}\) expands \(E_u\) and contracts \(E_s\) with a
spectral gap.  If the nonlinear part is Lipschitz-small in a neighborhood of
the fixed point, then there is a local stable manifold tangent to \(E_s\).
Points on this manifold remain in the neighborhood under forward RG
iteration and converge to \(V_\ast\).
\end{theorem}

\begin{proof}
Translate \(V_\ast\) to the origin and write \(V=(u,s)\) according to
\(E_u\oplus E_s\).  The RG map has the form
\[
  (u,s)\longmapsto (Au+F_u(u,s),Bs+F_s(u,s)),
\]
where \(A^{-1}\) and \(B\) are contractions and \(F=(F_u,F_s)\) has small
Lipschitz constant on a ball.  For a candidate graph \(u=h(s)\), invariance
requires
\[
  h(Bs+F_s(h(s),s))=Ah(s)+F_u(h(s),s).
\]
This equation defines a contraction on the complete metric space of Lipschitz
graphs with \(h(0)=0\) and \(Dh(0)=0\), after the ball is chosen so that the
nonlinear Lipschitz constant is dominated by the spectral gap.  The fixed
graph is invariant.  Iterating the contraction estimates gives convergence
to the fixed point for initial data on the graph.
\end{proof}

\section{Polymer Representation}

On a lattice, a polymer is a finite union of blocks.  A polymer activity
\[
  K(X,\phi)
\]
assigns a functional to each polymer \(X\).  The partition function is written
as
\[
  Z=\int d\mu_C(\phi)\,e^{-V_{\mathrm{loc}}(\phi)}
  \sum_{\mathcal X}
  \prod_{X\in\mathcal X}K(X,\phi),
\]
where \(\mathcal X\) ranges over compatible families of disjoint polymers.
The RG step integrates fluctuation fields inside each block, extracts local
Taylor coefficients into \(V_{\mathrm{loc}}\), and maps the remaining
nonlocal dependence into a new activity \(K'\).  The proof obligation is an
estimate of the form
\[
  \|K'\|_{\mathrm{polymer}}
  \leq c L^{-\alpha}\|K\|_{\mathrm{polymer}}
  +c'\|K\|_{\mathrm{polymer}}^2
\]
with declared constants and norm.  This is the mechanism by which irrelevant
remainders are controlled in constructive RG.

\section{Fixed Points And Continuum Limits}

A continuum limit is obtained by tuning the unstable coordinates so that the
RG trajectory remains near a fixed point for arbitrarily many steps and then
converges, after reconstruction, to regulator-independent correlation
functions.  The existence of a non-Gaussian fixed point as a Banach-space
fixed point is a theorem only after the RG map, norm, domain, and estimates
are supplied.  Perturbative beta functions and numerical fixed-point searches
are evidence for a coordinate chart; they do not replace the fixed-point
construction.

\section{Relation To Universality Classes}

In this monograph, a universality class is a triple
\[
  (\mathcal R,\mathcal N,V_\ast)
\]
together with an equivalence relation on microscopic regulators whose tuned
RG trajectories converge to \(V_\ast\) in the normed neighborhood
\(\mathcal N\) and whose reconstructed long-distance observables agree after
declared normalizations.  The phrase therefore denotes a constructed
attraction statement, not a label assigned solely by symmetry and dimension.

\begin{openproblem}[Nonperturbative Wilsonian fixed points]
\label{op:constructive-rg-nonperturbative-fixed-points}
Construct non-Gaussian Wilsonian fixed points in three and four dimensions as
fixed points of explicit RG maps on Banach spaces of interactions, with
stable manifolds, operator spectra, and reconstruction of the associated
local QFT data.
\end{openproblem}
